
%\section{Naive Approach}


\begin{frame}{Refinement}
	True peak location will always be between 2 audio samples. 
	\begin{figure}
		\centering
		\includegraphics[width=0.7 \textwidth]{../../python_correlation_sandbox/modern/saved_figs/15_realistic_correlation_surface_around_expected_peak_10_of_block_period.pdf}
		\caption{Results zoomed and centered around true peak location}
	\end{figure}
\end{frame}

\begin{frame}{Refinement}
	\begin{align}
		(y_1 \star y_2) (\tau) & = T \sum_{m = -N/2}^{N/2-1} \overline{A_m} B_{m} \exp{(\twopii f_m \tau)}
	\end{align}
	\vspace{0.5cm}
	\begin{align}
		\diff{^{(n)}}{\tau^{(n)}} (y_1 \star y_2)(\tau) & = (\twopii)^{n} T \sum_{m = -N/2}^{N/2-1} f_m^n \overline{A_m} B_m \exp{(\twopii f_m \tau)} \\
		&\text{where $f^n_{-N/2} = 0$ for odd n}
	\end{align}
\end{frame}


%\begin{frame}{xxxxxxx}
%	\begin{figure}
%		\centering
%		\includegraphics[width=0.7 \textwidth]{../../python_correlation_sandbox/modern/saved_figs/12_realistic_correlation_surface_around_expected_peak_500_of_block_period.pdf}
%		\caption{Correlation surface}
%	\end{figure}
%\end{frame}


\begin{frame}{Correlation Surface Derivatives}
	The correlation peak can be found from the zero-crossing of the first derivative. Let's put it in a root solver and let it run a couple iterations. 
	
	
	\centering
%	\begin{align*}
%		sinc(t) &= \frac{sin(t)}{t} \\
%		\mathcal{F}\{sinc(t)\} &= rect(\omega)
%	\end{align*}
	\begin{columns}
		\begin{column}{0.5\textwidth}
			\centering
			\includegraphics[width=\linewidth]{../../python_correlation_sandbox/modern/saved_figs/15_realistic_correlation_surface_around_expected_peak_10_of_block_period.pdf}
			%			\captionof{figure}{$sinc(t)$}
		\end{column}
		\begin{column}{0.5\textwidth}
			\centering
			\includegraphics[width=\linewidth]{../../python_correlation_sandbox/modern/saved_figs/16_realistic_fracddtau_correlation_surface_around_expected_peak_10_of_block_period.pdf}
			%			\captionof{figure}{Second Caption}
		\end{column}
	\end{columns}
\end{frame}


\begin{frame}[fragile]{Newton Solver}
	Solving for the corr peak requires evaling the 0th, 1st, and 2nd derivatives
		\begin{align}
			x_{i+1} \leftarrow x_i - \frac{f'(x_i)}{f''(x_i)}
		\end{align}
	
	
	\begin{lstlisting}[basicstyle=\ttfamily\scriptsize, 
		linebackgroundcolor={%
			\ifnum\value{lstnumber}=9 \color{yellow!30}\fi
			\ifnum\value{lstnumber}=10 \color{yellow!30}\fi
			\ifnum\value{lstnumber}=11\color{yellow!30}\fi
		}]
struct NewtonSolver
{
	template <typename CallableT>
	static Solution solve(const CallableT &fd0_fd1_fd2, Precision guess_tau_s, size_t n_steps = 4)
	{
		Precision x_n = guess_tau_s;
		for (size_t i = 0; i < n_steps; ++i)
		{
			auto [f_d0, f_d1, f_d2] = fd0_fd1_fd2(x_n);
			auto update = f_d1 / f_d2;
			x_n = x_n - update;
		}
		
		auto [f_d0, f_d1, f_d2] = fd0_fd1_fd2(x_n);
		return Solution{.x_final=x_n, .f_at_x_final=f_d0};
	}
};
	\end{lstlisting}
\end{frame}

\begin{frame}{The Refinement Problem}
	\vspace{-0.5cm}
	Fundamentally different problem now since we need only data point, but at a completely arbitrary location
	\begin{align*}
		\begin{bNiceMatrix}
			\\
			\\
			x[12.34]\\
			\\
			\\
		\end{bNiceMatrix} &=
		\begin{bmatrix}
			&  &  &  & \\
			&  &  &  & \\
			1 & e^{i\frac{2\pi(12.34)}{N}} & e^{i\frac{4\pi(12.34)}{N}} & \cdots & e^{i\frac{2\pi(12.34)(N-1)}{N}}\\
			 &  &  &  & \\
			 & & & & \\
			 & & & & % \\
		\end{bmatrix}
		\begin{bNiceMatrix}
			X[0]\\
			X[1]\\
			X[2]\\
			\vdots\\
			X[N-1]
		\end{bNiceMatrix}
	\end{align*}
\end{frame}


\begin{frame}{Refinement Approach 1}
	\begin{itemize}
		\item One function to compute the function at any derivative order
		\item Evaluate multiple derivatives simultaneously via template metaprogramming
		\item Compiler cannot optimize or reuse calls to std::exp or std::polar
	\end{itemize}
\end{frame}

\begin{frame}[fragile]{Refinement Approach 1}
	\begin{lstlisting}[basicstyle=\ttfamily\scriptsize, 
		linebackgroundcolor={%
			\ifnum\value{lstnumber}=1 \color{yellow!30}\fi
			\ifnum\value{lstnumber}=7 \color{yellow!30}\fi
			\ifnum\value{lstnumber}=8 \color{yellow!30}\fi
			\ifnum\value{lstnumber}=24 \color{yellow!30}\fi
		}]
template <size_t deriv_order, typename TauT>
auto correlate_impl(const CoeffsT &B, const TauT tau) const
{
	Real sum(0.0);
	for (size_t i = 0; i < Ncoeffs; ++i)
	{
		auto Wn = std::polar(1, constants::twopi * freqs[i] * tau);
		Complex term_i = coeffs_a_conj_and_f[deriv_order][i]*B[i]*Wn;
		
		// pull only the real/imag parts. components not taken are 0
		if constexpr(deriv_order%4==0)      sum+=term_i.real();
		else if constexpr(deriv_order%4==1) sum-=term_i.imag();
		else if constexpr(deriv_order%4==2) sum-=term_i.real();
		else if constexpr(deriv_order%4==3) sum+=term_i.imag();
	}
	return sum;
}
	\end{lstlisting}
	
\end{frame}


\begin{frame}[fragile]{Refinement Approach 1}
	\begin{lstlisting}[basicstyle=\ttfamily\scriptsize, 
		linebackgroundcolor={%
			\ifnum\value{lstnumber}=6 \color{yellow!30}\fi
			\ifnum\value{lstnumber}=14 \color{yellow!30}\fi
		}]
template <typename TauT, size_t... indicesT>
auto correlate_and_derive_impl(const CoeffsT &B, const TauT tau, std::integer_sequence<size_t, indicesT...>) const
{
	using correlation_type = decltype(correlate_impl<0>(B, tau));
	std::array<correlation_type, sizeof...(indicesT)> ret;
	((ret[indicesT] = correlate_impl<indicesT>(B, tau)), ...);
	return ret;
}


template <size_t deriv_order, typename TauT = Real>
auto correlate_and_derive_v0(const CoeffsT &B, const TauT tau) const
{
	return correlate_and_derive_impl(B, tau, std::make_index_sequence<deriv_order + 1>());
}
	\end{lstlisting}
	
\end{frame}

\begin{frame}{Results}
	\refinementtable{3}{1}
\end{frame}


\begin{frame}{Refinement Approach 2}
	\begin{itemize}
		\item Compute all derivative orders in the same function
		\item Evaluate all twiddles on each iteration
	\end{itemize}
\end{frame}

\begin{frame}[fragile]{Refinement Approach 2}
	\begin{lstlisting}[basicstyle=\ttfamily\scriptsize, 
		linebackgroundcolor={%
			\ifnum\value{lstnumber}=8 \color{yellow!30}\fi
			\ifnum\value{lstnumber}=9 \color{yellow!30}\fi
			\ifnum\value{lstnumber}=13 \color{yellow!30}\fi
		}]
template <size_t deriv_order, typename TauT = Real>
auto correlate_and_derive_v1(const CoeffsT &B, const TauT tau) {
	constexpr size_t Norders = deriv_order + 1; // include orig fcn
	
	std::array<Real, Norders> sum; // fill 0
	for (size_t i = 0; i < Ncoeffs; ++i)
	{
		auto Wn = std::polar(1, twopi * freqs[i] * tau);
		auto coeff_b_and_Wn = B[i] * Wn;
		
		for (size_t order = 0; order < Norders; ++order)
		{
			Complex term_i=coeffs_a_conj_and_f[order][i]*coeff_b_and_Wn;
			
			// pull only the real/imag parts. components not taken are 0
			if constexpr(derivative_order%4==0)      sum+=term_i.real();
			else if constexpr(derivative_order%4==1) sum-=term_i.imag();
			else if constexpr(derivative_order%4==2) sum-=term_i.real();
			else if constexpr(derivative_order%4==3) sum+=term_i.imag();
		}
	} return sum;
}
	\end{lstlisting}
	
\end{frame}


\begin{frame}{Results}
	\refinementtable{3}{2}
\end{frame}


\begin{frame}{Discrete Fourier Transform (Summation Form)}
	\begin{align}
		x[k] &= \sum_{n=0}^{N-1} X[n] e^{i\frac{2\pi}{N}kn}, \qquad k=0,1,\ldots,N-1
	\end{align}
	\vspace{0.5cm}
	\begin{align}
	\tikzmarknode{precompute}{W} &= e^{i\frac{2\pi}{N}\tau} \\
	x(k) &= \sum_{n=0}^{N-1} X[n] \tikzmarknode{update}{W^n}, \qquad k=0,1,\ldots,N-1
	\end{align}
	
	\begin{tikzpicture}[overlay, remember picture]
		% Draw a diagonal arrow/line from the term to an offset coordinate and place text
		\draw[<-, thick, red] (precompute.west) -- ++(-1.0cm, 0.5cm) 
		node[annot, above] {Precompute once};
	\end{tikzpicture}
	
	\begin{tikzpicture}[overlay, remember picture]
	% Draw a diagonal arrow/line from the term to an offset coordinate and place text
	\draw[<-, thick, red] (update.south) -- ++(0.5cm, -1.0cm) 
	node[annot, below] {Update iteratively since $W^n = W^{n-1} \ W$};
	\end{tikzpicture}
\end{frame}

\begin{frame}{Refinement Approach 3}
	\begin{itemize}
		\item Compute all derivatives in the same function
		\item Evaluate one twiddle. Use it to recursively compute the rest
		\item Potential accuracy hit for some scenarios 
	\end{itemize}
\end{frame}


\begin{frame}[fragile]{Refinement Approach 3}
	\begin{lstlisting}[basicstyle=\ttfamily\scriptsize, 
		linebackgroundcolor={%
			\ifnum\value{lstnumber}=5 \color{yellow!30}\fi
			\ifnum\value{lstnumber}=6 \color{yellow!30}\fi
			\ifnum\value{lstnumber}=15 \color{yellow!30}\fi
			\ifnum\value{lstnumber}=22 \color{yellow!30}\fi
		}]
template <size_t deriv_order, typename TauT = Real>
auto correlate_and_derive_v2(const CoeffsT &B, const TauT tau) {
	constexpr size_t Norders = deriv_order + 1; // include orig fcn
	
	Complex Wn(1,0);
	Complex W1 = std::polar(1, twopi * freqs[1] * tau);
	
	std::array<Real, Norders> sum; // fill 0
	for (size_t i = 0; i < Ncoeffs; ++i)
	{
		auto coeff_b_and_Wn = B[i] * Wn;
		
		for (size_t order = 0; order < Norders; ++order)
		{
			Complex term_i=coeffs_a_conj_and_f[order][i]*coeff_b_and_Wn;
			
			if (order % 4 == 0)          sum[order] += term_i.real();
			else if (order % 4 == 1)     sum[order] -= term_i.imag();
			else if (order % 4 == 2)     sum[order] -= term_i.real();
			else /*if (order % 4 == 3)*/ sum[order] += term_i.imag();
		}
		Wn *= W1;
	} return sum;
}
	\end{lstlisting}
	
\end{frame}


\begin{frame}{Results}
	\refinementtable{3}{3}
\end{frame}




